\documentclass{article}
\usepackage[a4paper, left=1.5cm, right=1.5cm, top=1cm, bottom=2cm]{geometry}

\usepackage{boldline}
\usepackage{tikz,tcolorbox}
\usepackage{amsmath}
\usepackage[table,xcdraw]{xcolor}
\usepackage{listings}
\usepackage{array,multirow} % For customizing tables
\usepackage{booktabs} % For better horizontal lines
\usepackage{makecell}
\setlength{\parindent}{0pt}
\usepackage{siunitx}
\usepackage{tkz-tab}
\usepackage{amssymb}
\usepackage{amsmath}
\usepackage{enumitem}

\usepackage{caption}
\usepackage{float}


\newcommand{\exer}[1]{
  \section*{Exercice #1}
  \vspace{-0.5cm}
  \noindent\rule{\textwidth}{0.5pt}%
}

\newcommand{\tit}[1]{
\begin{center}
    \Large{\textbf{{#1}}}
\end{center}
}

\definecolor{commentgray}{HTML}{676160}
\definecolor{messagegreen}{HTML}{17B867}
\definecolor{myblue}{HTML}{10C2C4}

\tcbuselibrary{skins, breakable, theorems}
\usepackage{booktabs} 

\newtcolorbox{prettyBox}[2]{
  enhanced,
  colback=white!90!#2,   % Background color based on the second parameter (color)
  colframe=#2!60!black,  % Frame color based on the second parameter (color)
  coltitle=white,        % Title color (white)
  fonttitle=\bfseries\Large,
  title=#1,              % Title from the first parameter
  boxrule=1mm,
  arc=0.5mm,
  drop shadow=#2!35!gray, % Drop shadow color based on the second parameter (color)
}



\lstdefinestyle{cmd}{
 basicstyle=\ttfamily,
 backgroundcolor=\color{lightgray!20},
 frame=single
}

\lstdefinestyle{pythonstyle}{
    language=python,                    % Language set to Python
    basicstyle=\ttfamily\footnotesize,   % Change basic font size
    keywordstyle=\color{blue}\bfseries, % Different keyword style
    stringstyle=\color{red},         % Different string color
    commentstyle=\color{green!60!black}\itshape, % Adjust comment color
    numbers=left,                       % Line numbers on the left
    numberstyle=\tiny\color{gray},      % Smaller number font and color
    stepnumber=1,                       % Number each line
    frame=single,                       % Single frame around code
    tabsize=4,                          % Adjust tab size
    showstringspaces=false,             % Do not show spaces in strings
    captionpos=b,% Position of caption
    breaklines=true,
    inputencoding=utf8
}




\begin{document}
\tit{TD N\(^{\boldsymbol{\circ}}\)\hspace{0.1cm}3}
\exer{1}

\vspace{0.25cm}


\begin{prettyBox}{Naive Bayes}{myblue}
    Naive Bayes est un algorithme supervisé de classification basé sur le théorème de 
    Bayes naïf et les probabilités conditionnelles. Il prend en entrée un point $X$ et donne
    sa classe $Y$.
\end{prettyBox}

\vspace{0.25cm}


\begin{prettyBox}{Formule}{myblue}
\[ P(Y|X) = \dfrac{P(X|Y) \cdot P(Y)}{P(X)} \] 

\begin{itemize}
    \setlength{\itemsep}{0pt}
    \item $X$ : un point.
    \item $Y$ : une classe.
    \item $P(Y|X)$ : Probabilité postérieure.
    \item $P(X|Y)$ : Probabilité conditionnelle (likelihood).
    \item $P(Y)$ : Probabilité priorie.
    \item $P(X)$ : Probabilité d'évidence.
\end{itemize}
\end{prettyBox}

\vspace{0.25cm}

\begin{prettyBox}{Algo}{myblue}
\begin{enumerate}
    \setlength{\itemsep}{0pt}
    \item Calculer les probabilités priors de toutes les classes $Y$.
    \item Calculer les probabilités conditionnelles de toutes les classes $Y$ :
        \[ P(X|Y) = P(x_1 \cap \dots \cap x_M | Y) = \prod_{i=1}^M P(x_i|Y) \]
        \begin{itemize}
            \item $M$ : nombre d'attributs
            \item $x_i$ : i-ème attribut
        \end{itemize}
        On suppose que les attributs sont indépendants les uns des autres.
    \item Calculer le numérateur de $P(Y|X)$ c'est-à-dire $P(X|Y) \cdot P(Y)$ pour toutes les classes $Y$ et prendre la classe avec la plus grande probabilité.
    \item Le plus grande probabilité déterminera la classe du point $X$.
    \item On n'a pas besoin de calculer le dénominateur $P(X)$, mais voici sa formule :
        \begin{align*}    
            P(X) &= P(X \cap C) \\[0.1cm]
                  &= P(X \cap (Y_1 \cup \dots \cup Y_j)) \\[0.1cm]
                  &= P(X\cap Y_1) + \dots + P(X\cap Y_j) \\[0.1cm]
                  &= P(X|Y_1) \cdot P(Y_1) + \dots + P(X|Y_j) \cdot P(Y_j)
        \end{align*} 
        \begin{itemize}
            \item $C$ : représente toutes les classes possibles.
            \item $j$ : nombre de classes.
            \item $Y_i$ : i-ème classe.
        \end{itemize}
\end{enumerate}
\end{prettyBox}

\newpage

\begin{table}[h]
\centering

\renewcommand{\arraystretch}{2.25} 
\setlength{\tabcolsep}{12pt}      

\begin{tabular}{|c|c|c|c|} 

\hline

Jour & Temps ($x_1$) & Vent($x_2$) & Jouer au Tennis ($C$) \\ \hline
1 & Pluie & Oui & Non \\ \hline
2 & Ensoleillé & Non & Oui \\ \hline
3 & Pluie & Oui & Oui \\ \hline
4 & Pluie & Non & Oui \\ \hline
5 & Ensoleillé & Oui & Non \\ \hline
6 & Pluie & Non & Oui \\ \hline
7 & Ensoleillé & Oui & Non \\ \hline
8 & Pluie & Non & Oui \\ \hline
9 & Ensoleillé & Non & Non \\ \hline
10 & Pluie & Oui & Oui \\ \hline
\end{tabular}
\end{table}

Utilisez le classifieur Naïve Bayes pour prédire si la personne jouera au tennis étant donné
les nouvelles conditions suivantes : $X_{\text{test}}$ = (Temps= Ensoleillé, Vent= Non) 

\vspace{0.75cm}
Calcule des probabilites priorie :
\[P(\text{Oui}) = \dfrac{6}{10} \quad , \quad P(\text{Non}) = \dfrac{4}{10}\]


\vspace{0.75cm}
Calcule des probabilites conditionelle :
\begin{align*}
    &P(X_{\text{test}} = P (\text{Ensoleillé , Non}) \hspace{0.1cm} | \hspace{0.1cm} Y = \text{Oui}) = P(x_1 = \text{Ensoleillé} \hspace{0.1cm} | \hspace{0.1cm} Y =\text{Oui}) \cdot P(x_2 = \text{Non} \hspace{0.1cm} | \hspace{0.1cm} Y =\text{Oui}) = \dfrac{1}{6} \cdot \dfrac{4}{6} = \dfrac{1}{9}\\[0.15cm] 
    &P(X_{\text{test}} = P (\text{Ensoleillé , Non}) \hspace{0.1cm} | \hspace{0.1cm} Y = \text{Non}) = P(x_1 = \text{Ensoleillé} \hspace{0.1cm} | \hspace{0.1cm} Y =\text{Non}) \cdot P(x_2 = \text{Non} \hspace{0.1cm} | \hspace{0.1cm} Y =\text{Non}) = \dfrac{3}{4} \cdot \dfrac{1}{4} = \dfrac{3}{16} 
\end{align*}

\vspace{0.75cm}
Calcule des numerateurs :

\begin{align*}
    &\text{Numerateur(Oui)} = P (\text{Ensoleillé , Non}) \hspace{0.1cm} | \hspace{0.1cm} Y = \text{Oui})\cdot P(\text{Oui}) = \dfrac{1}{9} \cdot \dfrac{6}{10} = 0.067\\[0.1cm]
    &\text{Numerateur(Non)} = P (\text{Ensoleillé , Non}) \hspace{0.1cm} | \hspace{0.1cm} Y = \text{Non})\cdot P(\text{Non}) = \dfrac{3}{16} \cdot \dfrac{4}{10} = 0.075
\end{align*}

\vspace{0.45cm}

Puisque Numerateur(Non) = 0.075 $>$ Numerateur(Oui) = 0.067 le classifieur Naïve Bayes prédit la classe Non.\\[0.1cm]
\textbf{Conclusion} : Selon les conditions de test ($X_{\text{text}}$ = (Temps= Ensoleillé, Vent= Non)) l'algorithme prédit que la personne ne jouera pas au tennis.\\[0.15cm]

Pour l'information la probabilité postérieure réelle serait :
\[P(Y = \text{Non} | X_{\text{test}}) = \dfrac{0.075}  {(0.067 +0.075)} = 0,53\] 

\end{document}
